\documentclass[12pt]{article}

\usepackage[utf8]{inputenc}
\usepackage[margin=1in]{geometry}
\usepackage{setspace}
\usepackage{url}
\usepackage{mathptmx}

\doublespacing
\setlength{\emergencystretch}{3em}
\setlength{\parindent}{0pt}
\setlength{\parskip}{0.25em}

\begin{document}

\begin{center}
\textbf{Comparison Essay Working Plan}
\end{center}

Ma Toan Bach\\
Student Number: 112708227\\
LSO 100--Canadian Short Story (KCC)\\
Date: July 31, 2026

\vspace{0.5em}
\textbf{Selected Stories}

Alice Munro's ``Boys and Girls'' and Madeleine Thien's ``Simple Recipes.''

\textbf{Working Title}

\textit{Family Pressure, Identity, and Belonging in ``Boys and Girls'' and ``Simple Recipes''}

\textbf{Working Thesis}

Both ``Boys and Girls'' and ``Simple Recipes'' show that family shapes identity through love, expectation, and control; however, Munro emphasizes the quiet normalization of gender inequality, while Thien shows how cultural conflict and violence fracture belonging from within the family.

\textbf{Essay Structure}

\textbf{Introduction}

Introduce both stories and authors. Briefly state that both stories use first-person narration to examine how family relationships shape a young narrator's sense of identity. End with the working thesis above or a revised version in your own words.

\textbf{Body Paragraph 1: Fathers, Authority, and Recognition}

Main point: Both stories begin with father figures who seem connected to care, skill, and family identity, but both fathers also become sources of limitation or fear.

Munro evidence to use: the father calling the narrator his ``new hired hand''; the narrator's pride in outdoor work; the later shift when her father accepts the idea that she is ``only a girl.''

Thien evidence to use: the rice-making ritual; the father teaching skill and tradition in the kitchen; the shift from care to violence at the dinner table.

Analysis goal: Explain that recognition from the father gives each narrator a sense of belonging, but that belonging becomes unstable when fatherly authority turns restrictive or harmful.

\textbf{Body Paragraph 2: Identity, Belonging, and Family Pressure}

Main point: Both narrators learn that family belonging can require them to accept roles or loyalties that conflict with their own understanding of themselves.

Munro evidence to use: pressure to move toward indoor work; the mother's and grandmother's expectations about proper girlhood; the Flora gate scene as resistance to obedience.

Thien evidence to use: the brother's rejection of the meal; the narrator's divided loyalty after witnessing violence; the narrator's fear of becoming complicit by continuing to love her father.

Analysis goal: Compare gender pressure in Munro with cultural and emotional pressure in Thien. Show that both stories treat family as a place of belonging and conflict at the same time.

\textbf{Body Paragraph 3: Memory and Retrospective Narration}

Main point: Both stories use adult reflection on childhood memory to show that identity is understood gradually, after painful family experiences have been reconsidered.

Munro evidence to use: the narrator's changing understanding of the word ``girl''; the ending phrase ``Maybe it was true.''

Thien evidence to use: the adult narrator's memory of the kitchen; the reflection on shame, grief, and the difficulty of reconciling love with fear.

Analysis goal: Explain the difference in tone. Munro's narration traces a gradual social awakening, while Thien's narration returns to trauma in order to understand divided love and damaged belonging.

\textbf{Conclusion}

Return to the thesis without repeating it word for word. Emphasize that the comparison shows how Canadian short fiction often treats family as both a source of identity and a source of pressure. End with a broader insight about belonging: home can shape the self, but it can also expose the limits placed on a person by gender, culture, authority, and memory.

\textbf{Quote Checklist Before Drafting}

Use at least two short quotations from each story. Keep quotations brief and analyze each one directly. Add page numbers from the exact editions or course copies you use.

\textbf{Secondary Source Plan}

Use Nischik to support the discussion of gender roles and the narrator's resistance in ``Boys and Girls.'' Use Ty to support cultural context for Asian Canadian literature, immigration, and belonging in relation to Thien. The essay should still rely mainly on close reading of the two stories.

\clearpage
\begin{center}
\textbf{References}
\end{center}
\raggedright

\hangindent=0.5in\noindent Munro, A. (1968). \textit{Dance of the happy shades}. Ryerson Press.

\hangindent=0.5in\noindent Nischik, R. M. (2007). (Un-)doing gender: Alice Munro, ``Boys and Girls'' (1964). In R. M. Nischik (Ed.), \textit{The Canadian short story: Interpretations} (pp. 203--218). Camden House. \url{https://doi.org/10.1515/9781571136886-016}

\hangindent=0.5in\noindent Thien, M. (2001). \textit{Simple recipes}. McClelland \& Stewart.

\hangindent=0.5in\noindent Ty, E. (2016). (East and Southeast) Asian Canadian literature: The strange and the familiar. In C. Sugars (Ed.), \textit{The Oxford handbook of Canadian literature} (pp. 564--582). Oxford University Press. \url{https://doi.org/10.1093/oxfordhb/9780199941865.013.31}

\end{document}
